\chapter{Introduction}

\label{intro}

\section{Context and Motivation}

Music performances are not experienced through sound alone. In live concerts, club environments, and electronic music events, lighting plays a central role in shaping atmosphere, directing attention, and amplifying emotional intensity. Changes in colour, brightness, movement, and timing can make a musical moment feel larger, more intimate, more aggressive, or more immersive. For electronic music in particular, where repetition, gradual build-ups, drops, and timbral changes are essential expressive devices, visual response becomes part of the perceived structure of the performance.

At the same time, expressive stage lighting is difficult to automate. Traditional lighting design depends on human expertise, preparation time, and detailed manual programming. Automatic Stage Lighting Control addresses this problem by generating lighting parameters from music, with the potential to support live performances, influence the audience's emotional experience, and provide a more accessible alternative to dedicated lighting engineers. However, many prior approaches rely on a two-stage strategy: first classifying music into predefined style or emotion categories, then assigning lighting patterns to those categories. Such pipelines can be useful, but they also reduce the richness of music to a limited set of labels before the visual output is generated.

Recent generative approaches suggest a more flexible direction. Instead of mapping a classified category to a fixed lighting rule, Skip-BART treats automatic stage lighting as a sequence generation problem and learns to produce lighting sequences directly from music. This framing is especially relevant for creative applications because lighting design is not only a technical control problem, but also an artistic process involving timing, continuity, contrast, and expressive variation. A system that can use such models in real time must therefore solve more than an inference problem. It must capture or play audio, extract and stream meaningful signals, run analysis and generative models, transport outputs with low overhead, and update an interactive visual scene quickly enough for the result to feel synchronized with the music.

The motivation for this thesis stems from the convergence of two passions: programming and electronic dance music. Programming offers the means to build systems, structure complex processes, and turn abstract ideas into working applications. Electronic dance music provides a creative world where rhythm, energy, and atmosphere are inseparable from the visual experience — and where the connection between sound and light is felt immediately on the dance floor. This thesis was the natural opportunity to bring both together in a single project, exploring how music analysis and generative lighting models can be integrated into an application that is not only technically functional, but also visually engaging and suitable for interactive demonstration.

\section{Objectives}

The objective of the thesis is to design and implement a real-time sound-to-light prototype for generative stage lighting. The system analyzes audio while it is playing, extracts information relevant to musical timing and visual control, and uses model outputs to update a virtual stage environment. Beyond the technical pipeline, the goal is to create an environment in which a user can freely experiment — testing scene ideas, exploring lighting concepts, and interacting with the visual output as a creative tool rather than a static visualization. The full source code of the prototype is publicly available at \url{https://github.com/alexapvl/thesis_project}.

The application follows a hybrid architecture: the browser handles audio input, playback, user interaction, and 3D scene rendering, while a local server performs beat detection and generative lighting inference. Audio data is streamed from the browser to the server, and musical and lighting parameters are sent back to drive real-time visual updates.

A central design concern is real-time responsiveness, and it means different things for the two kinds of visual output the system produces. The thesis distinguishes a \emph{beat-synchronization path}, in which detected beats drive immediate visual reactions, from a \emph{generative lighting path}, in which Skip-BART produces the evolving colour palette. The system does not need to satisfy the extremely strict action-to-sound constraints of a musical instrument, but the beat-synchronization path must still keep visual reactions close enough to musical events to be perceived as synchronized; for this path the architecture targets a sub-100 ms end-to-end threshold and estimates it at approximately 90 ms. The generative lighting path operates on a coarser timescale by design: because Skip-BART generates lighting for a whole audio window at once rather than streaming frame by frame, the colour palette refreshes every few seconds rather than within milliseconds. It is therefore governed by an update \emph{cadence} rather than the sub-100 ms threshold---a distinction developed in \fullref{chap:methodology-requirements-system-design}{Chapter} and quantified in \fullref{chap:evaluation}{Chapter}. For this reason, the thesis treats system design as the management of latency, modularity, and usability under the constraints of the chosen platform.

\section{Contributions}

The original contribution of this thesis is the design, implementation, and evaluation of a complete real-time sound-to-light prototype that integrates beat detection and a generative lighting model into a single interactive application. The work does not propose a new machine-learning model; instead, its novelty lies in the system-level engineering required to make an offline generative model usable in a live, browser-based setting. The main contributions are the following:

\begin{itemize}
    \item A \textbf{hybrid browser--server architecture} that separates real-time audio capture, user interaction, and 3D rendering in the browser from beat detection and generative lighting inference on a local server, connected through a low-overhead streaming protocol.
    \item The \textbf{integration of Skip-BART, an offline generative lighting model, into a live streaming pipeline}, including the windowing and cadence strategy required to drive a continuously evolving colour palette from a model that generates lighting one audio window at a time.
    \item A \textbf{dual-path real-time design} that explicitly separates a low-latency path for beat synchronization from a coarser-grained generative lighting path, each governed by its own responsiveness target.
    \item A \textbf{typed communication protocol} at the browser--server boundary that structures the exchange of audio data, musical parameters, and lighting outputs.
    \item An \textbf{empirical evaluation} of the prototype, benchmarking end-to-end latency, server-side processing, and inference cost across two hardware configurations.
\end{itemize}

The original results of the thesis are therefore primarily systemic and quantitative: a working prototype that demonstrates the feasibility of real-time generative stage lighting on the web, together with latency and performance measurements that characterise its behaviour. These results are presented in detail in \fullref{chap:evaluation}{Chapter}.

\section{Use of Generative AI Tools}

In the interest of transparency, the author declares that generative AI tools were used in a limited and supporting capacity during the preparation of this thesis. Specifically, Claude (Anthropic) was used to improve the readability of the text, to rephrase passages that had already been written by the author, and to assist with selected aspects of software development, such as debugging, refactoring, and drafting small code fragments that were subsequently reviewed, adapted, and integrated by the author. These tools were not used to generate the original ideas, the system architecture, the research contributions, the experimental results, or the conclusions of this work, all of which are the author's own. Every AI-assisted suggestion, whether textual or programmatic, was reviewed, edited, and validated by the author, who takes full responsibility for the entire content of the thesis, for the source code of the prototype, and for the results obtained. Generative AI tools are not credited as authors, are not treated as reliable sources of factual content, and do not replace the author's own intellectual contribution.

\section{Thesis Structure}

The thesis is organized as follows. \fullref{chap:theoretical-background}{Chapter} introduces the theoretical foundations of the work, covering digital audio and Music Information Retrieval, rhythmic structure, beat tracking, music-visual mapping, automated stage lighting control, real-time processing, and the deep-learning architectures relevant to the system. \fullref{chap:technologies-and-platform-constraints}{Chapter} evaluates the web as a platform for real-time audio-visual applications, examining rendering with WebGL and Three.js, audio processing through the Web Audio API, browser-based machine-learning deployment, communication mechanisms, frontend framework selection, and the design implications of platform constraints. \fullref{chap:methodology-requirements-system-design}{Chapter} translates those constraints into a concrete system design by defining the methodological approach, requirements, hybrid architecture, module responsibilities, frontend data flow, and the end-to-end runtime pipeline. \fullref{chap:implementation}{Chapter} describes how that design is realised in code, focusing on the engineering problems that most shaped the prototype: the typed protocol at the browser--server boundary, the integration of an offline generative model into a live streaming pipeline, the real-time beat path, and the separation between runtime and editable scene state. \fullref{chap:evaluation}{Chapter} evaluates the prototype empirically, benchmarking end-to-end latency, server-side processing, and inference cost across two hardware configurations. \fullref{conclusions}{Chapter} concludes the thesis by consolidating the outcomes of the prototype and identifying the main lessons, limitations, and possible directions for future development.
